\section{PHÂN TÍCH YÊU CẦU HỆ THỐNG}

Trong phần này, nhóm thực hiện đi sâu vào việc phân tích các nhu cầu thực tế của người dùng đối với một hệ thống nhà thông minh đa tầng, từ đó chuyển hóa thành các yêu cầu chức năng và phi chức năng cụ thể cho hệ thống Mobile App và phần cứng Yolo:bit.

%=============================================================
\subsection{Yêu cầu người dùng (User Requirements)}
%=============================================================

Yêu cầu người dùng tập trung vào những gì người sở hữu ngôi nhà mong muốn đạt được khi sử dụng ứng dụng. Qua khảo sát thực tế, nhóm xác định các nhu cầu trọng tâm sau:

\begin{itemize}
    \item \textbf{UR1 - Quản lý không gian sống linh hoạt:} Người dùng cần một giao diện cho phép họ sơ đồ hóa ngôi nhà của mình (chia theo Tầng và Phòng) để dễ dàng định vị thiết bị thay vì quản lý một danh sách dài rời rạc.
    \item \textbf{UR2 - Khả năng giám sát từ xa:} Người dùng mong muốn biết được trạng thái nhiệt độ, độ ẩm và ánh sáng tại từng phòng trong nhà ở bất kỳ đâu thông qua điện thoại di động.
    \item \textbf{UR3 - Sự tiện nghi và tiết kiệm năng lượng:}
    \begin{itemize}
        \item Có thể bật/tắt quạt hoặc đèn mà không cần di chuyển đến công tắc vật lý.
        \item Hệ thống tự động tắt thiết bị khi không cần thiết (dựa trên cảm biến hoặc lịch trình) để tiết kiệm điện.
    \end{itemize}
    \item \textbf{UR4 - Tính an toàn và cảnh báo:} Người dùng cần được thông báo ngay lập tức nếu có các chỉ số môi trường bất thường (ví dụ: nhiệt độ quá cao có nguy cơ hỏa hoạn, ánh sáng ngoài ngưỡng cho phép).
    \item \textbf{UR5 - Tính minh bạch thông tin:} Người dùng muốn xem lại lịch sử hoạt động của các thiết bị để biết ai đã điều khiển hoặc hệ thống đã tự động kích hoạt vào lúc nào.
\end{itemize}

%=============================================================
\subsection{Yêu cầu chức năng (Functional Requirements)}
%=============================================================

Dựa trên yêu cầu người dùng, hệ thống được thiết kế để thực hiện các chức năng sau. Mỗi yêu cầu chức năng được mã hóa theo định dạng \textbf{FR-XX} để thuận tiện cho việc truy vết và kiểm thử.

\begin{table}[H]
\centering
\caption{Danh sách yêu cầu chức năng hệ thống}
\renewcommand{\arraystretch}{1.4}
\begin{tabular}{|m{2.2cm}|m{4cm}|m{7.5cm}|}
\hline
\textbf{Mã yêu cầu} & \textbf{Tên chức năng} & \textbf{Mô tả chi tiết} \\ \hline

FR-01 & Quản lý cấu trúc nhà & Cho phép người dùng khởi tạo và chỉnh sửa cấu trúc cây \textit{Nhà $\rightarrow$ Tầng $\rightarrow$ Phòng $\rightarrow$ Thiết bị}. Hỗ trợ thêm, sửa, xóa từng cấp. Khi một tầng/phòng bị xóa, tất cả thiết bị con phải được xử lý tương ứng trong cơ sở dữ liệu. \\ \hline

FR-02 & Điều khiển thiết bị & Hỗ trợ hai chế độ điều khiển: (1)~\textbf{Thủ công} -- người dùng trực tiếp bật/tắt quạt, chọn màu sắc và độ sáng đèn RGB qua giao diện ứng dụng; (2)~\textbf{Tự động} -- hệ thống tự điều chỉnh thiết bị dựa trên giá trị cảm biến so sánh với ngưỡng đã thiết lập. \\ \hline

FR-03 & Thiết lập lịch trình & Cho phép người dùng cài đặt kịch bản bật/tắt thiết bị theo thời gian (giờ/phút) và các ngày lặp lại trong tuần. Lịch trình được lưu vào database và thực thi ngay cả khi người dùng thoát ứng dụng. \\ \hline

FR-04 & Cảnh báo vượt ngưỡng & Người dùng thiết lập ngưỡng an toàn cho từng cảm biến. Khi giá trị thực tế vi phạm ngưỡng liên tục qua 3 chu kỳ gửi dữ liệu, hệ thống xác nhận trạng thái cảnh báo và gửi thông báo đẩy (push notification) tức thời đến điện thoại người dùng. \\ \hline

FR-05 & Đồng bộ dữ liệu MQTT & Hệ thống duy trì kết nối liên tục với Adafruit IO qua giao thức MQTT. Dữ liệu cảm biến từ mạch Yolo:bit được tự động ánh xạ đúng với từng phòng trên ứng dụng và cập nhật theo chu kỳ không quá 5 giây. \\ \hline

FR-06 & Giám sát thời gian thực & Hiển thị biểu đồ trực quan giá trị nhiệt độ, độ ẩm và ánh sáng theo thời gian thực cho từng phòng. Dữ liệu được cập nhật với độ trễ không quá 3 giây. \\ \hline

FR-07 & Nhật ký và báo cáo & Ghi lại toàn bộ thao tác người dùng, lệnh tự động và cảnh báo vào system log với các trường: Thời gian, Thiết bị, Hành động, Nguồn kích hoạt (User/Auto). Hỗ trợ lọc theo thời gian/phòng và xuất file báo cáo định kỳ. \\ \hline

\end{tabular}
\end{table}

%=============================================================
\subsection{Yêu cầu phi chức năng (Non-functional Requirements)}
%=============================================================

Để hệ thống vận hành trơn tru trong môi trường thực tế, bao gồm cả các tình huống mất kết nối hoặc tải cao, các chỉ tiêu kỹ thuật sau cần được đảm bảo:

\begin{table}[H]
\centering
\caption{Bảng đặc tả các yêu cầu phi chức năng}
\renewcommand{\arraystretch}{1.4}
\begin{tabular}{|m{3.2cm}|m{1.8cm}|m{1.8cm}|m{6.5cm}|}
\hline
\textbf{Nhóm} & \textbf{Mã chỉ số} & \textbf{Đặc tính} & \textbf{Giá trị / Chi tiết yêu cầu} \\ \hline

\multirow{2}{*}{\textbf{Hiệu năng}} 
    & NFR-P1 & Tốc độ điều khiển     & Thời gian phản hồi lệnh điều khiển từ App đến thiết bị $\leq$ 2 giây. \\ \cline{2-4}
    & NFR-P2 & Tốc độ cập nhật cảm biến & Giá trị cảm biến được cập nhật lên App với chu kỳ $\leq$ 5 giây/lần. \\ \hline

\multirow{2}{*}{\textbf{Độ tin cậy}} 
    & NFR-R1 & Uptime hệ thống        & Hệ thống hoạt động ổn định $\geq$ 95\% thời gian; server và thiết bị duy trì kết nối liên tục. \\ \cline{2-4}
    & NFR-R2 & Reconnect MQTT         & Khi kết nối MQTT bị ngắt, hệ thống tự động thử kết nối lại trong vòng 10 giây. Tỷ lệ thành công gửi/nhận dữ liệu $\geq$ 99\% trong điều kiện mạng ổn định. \\ \hline

\multirow{1}{*}{\textbf{Khả năng mở rộng}} 
    & NFR-S1 & Scalability            & Hệ thống quản lý đồng thời ít nhất 20 phòng mỗi tài khoản, mỗi phòng có đầy đủ cảm biến và thiết bị đầu ra độc lập. \\ \hline

\multirow{1}{*}{\textbf{Tính kịp thời}} 
    & NFR-T1 & Alert Latency          & Thông báo cảnh báo gửi đến App $\leq$ 3 giây kể từ khi cảm biến xác nhận trạng thái vi phạm ngưỡng. \\ \hline

\multirow{1}{*}{\textbf{Tính chính xác}} 
    & NFR-A1 & Scheduling Accuracy    & Sai số thực thi lịch trình $\leq$ 5 giây so với thời gian cài đặt. \\ \hline

\multirow{1}{*}{\textbf{Hoạt động offline}} 
    & NFR-O1 & Offline Resilience     & Khi mất kết nối Internet, vi điều khiển vẫn tự động điều khiển thiết bị dựa trên ngưỡng cảm biến đã được lưu cục bộ gần nhất. \\ \hline

\multirow{1}{*}{\textbf{Bảo mật}} 
    & NFR-SEC1 & Xác thực người dùng  & Chỉ tài khoản hợp lệ (đã đăng nhập) mới được truy cập hệ thống, thay đổi cấu trúc nhà hoặc thiết lập ngưỡng cảnh báo. \\ \hline

\end{tabular}
\end{table}

%=============================================================
\subsection{Mô hình Use-case hệ thống}
%=============================================================

\subsubsection{Danh sách các Use-case chính}

Hệ thống bao gồm các use-case chính sau, bao phủ toàn bộ luồng tương tác giữa người dùng, ứng dụng di động, server Adafruit IO và phần cứng Yolo:bit:

\begin{table}[H]
\centering
\caption{Danh sách use-case chính của hệ thống}
\renewcommand{\arraystretch}{1.4}
\begin{tabular}{|m{2cm}|m{4.5cm}|m{7.2cm}|}
\hline
\textbf{Mã UC} & \textbf{Tên Use-case} & \textbf{Mô tả ngắn} \\ \hline
UC-01 & Đăng nhập hệ thống         & Người dùng xác thực tài khoản để truy cập ứng dụng. \\ \hline
UC-02 & Quản lý cấu trúc nhà       & Thêm, sửa, xóa tầng/phòng và đăng ký thiết bị vào từng phòng. \\ \hline
UC-03 & Điều khiển thiết bị        & Bật/tắt và cấu hình quạt, đèn RGB ở chế độ thủ công hoặc tự động. \\ \hline
UC-04 & Thiết lập lịch trình       & Cài đặt thời gian tự động bật/tắt thiết bị theo lịch hàng tuần. \\ \hline
UC-05 & Thiết lập cảnh báo         & Cấu hình ngưỡng an toàn cho cảm biến và nhận thông báo khi vi phạm. \\ \hline
UC-06 & Nhận dữ liệu cảm biến      & Hệ thống tự động thu thập và hiển thị giá trị cảm biến theo thời gian thực. \\ \hline
UC-07 & Xem nhật ký và báo cáo     & Tra cứu lịch sử hoạt động, lọc theo thời gian/phòng và xuất báo cáo. \\ \hline
\end{tabular}
\end{table}

%-------------------------------------------------------------
\subsubsection{Đặc tả chi tiết các Use-case}
%-------------------------------------------------------------

% --- UC-01 ---
\paragraph{UC-01: Đăng nhập hệ thống}

\begin{table}[H]
\centering
\renewcommand{\arraystretch}{1.4}
\begin{tabular}{|p{3.8cm}|p{10cm}|}
\hline
\textbf{Tên Use-case}       & Đăng nhập hệ thống \\ \hline
\textbf{Mã Use-case}        & UC-01 \\ \hline
\textbf{Tác nhân}           & Người dùng \\ \hline
\textbf{Tiền điều kiện}     & Ứng dụng Mobile đã được cài đặt và có kết nối Internet. \\ \hline
\textbf{Luồng sự kiện chính} & 
1. Người dùng mở ứng dụng. \newline
2. Hệ thống hiển thị màn hình đăng nhập. \newline
3. Người dùng nhập tên tài khoản và mật khẩu. \newline
4. Hệ thống xác thực thông tin với server. \newline
5. Đăng nhập thành công, hệ thống chuyển đến màn hình Dashboard. \\ \hline
\textbf{Luồng ngoại lệ}     & 
\textit{E1 -- Sai thông tin:} Hệ thống hiển thị thông báo lỗi, yêu cầu nhập lại. \newline
\textit{E2 -- Mất kết nối:} Hệ thống hiển thị thông báo lỗi mạng, giữ nguyên màn hình đăng nhập. \\ \hline
\textbf{Hậu điều kiện}      & Người dùng được xác thực và có quyền truy cập toàn bộ chức năng. \\ \hline
\end{tabular}
\end{table}

% --- UC-02 ---
\paragraph{UC-02: Quản lý cấu trúc nhà}

\begin{table}[H]
\centering
\renewcommand{\arraystretch}{1.4}
\begin{tabular}{|p{3.8cm}|p{10cm}|}
\hline
\textbf{Tên Use-case}       & Quản lý cấu trúc nhà \\ \hline
\textbf{Mã Use-case}        & UC-02 \\ \hline
\textbf{Tác nhân}           & Người dùng \\ \hline
\textbf{Tiền điều kiện}     & Người dùng đã đăng nhập thành công (UC-01). Ứng dụng có kết nối Internet. \\ \hline
\textbf{Luồng sự kiện chính} & 
1. Người dùng truy cập mục \textit{Quản lý hệ thống}. \newline
2. Hệ thống hiển thị cấu trúc cây: Nhà $\rightarrow$ Tầng $\rightarrow$ Phòng $\rightarrow$ Thiết bị. \newline
3. Người dùng chọn thao tác: Thêm, Sửa hoặc Xóa một cấp bất kỳ. \newline
4. Người dùng nhập thông tin và nhấn \textit{Lưu}. \newline
5. Hệ thống yêu cầu xác nhận, người dùng đồng ý. \newline
6. Hệ thống cập nhật database và ghi nhật ký thao tác. \\ \hline
\textbf{Luồng ngoại lệ}     & 
\textit{E1 -- Xóa tầng/phòng còn thiết bị:} Hệ thống cảnh báo và xử lý xóa cascade toàn bộ thiết bị con. \newline
\textit{E2 -- Tên bị trùng:} Hệ thống hiển thị thông báo lỗi, yêu cầu đặt lại tên. \\ \hline
\textbf{Hậu điều kiện}      & Cấu trúc nhà được cập nhật; thao tác được lưu vào system log. \\ \hline
\end{tabular}
\end{table}

% --- UC-03 ---
\paragraph{UC-03: Điều khiển thiết bị}

\begin{table}[H]
\centering
\renewcommand{\arraystretch}{1.4}
\begin{tabular}{|p{3.8cm}|p{10cm}|}
\hline
\textbf{Tên Use-case}       & Điều khiển thiết bị \\ \hline
\textbf{Mã Use-case}        & UC-03 \\ \hline
\textbf{Tác nhân}           & Người dùng, Vi điều khiển Yolo:bit, Cảm biến ánh sáng \\ \hline
\textbf{Tiền điều kiện}     & Người dùng đã đăng nhập. Thiết bị đèn/quạt đang trực tuyến. \\ \hline
\textbf{Luồng sự kiện chính} & 
1. Người dùng chọn phòng cần điều khiển. \newline
2. Người dùng chọn chế độ \textbf{Thủ công} hoặc \textbf{Tự động}. \newline
\quad \textit{Thủ công:} Người dùng kéo thanh trượt độ sáng (0--100\%) và chọn màu sắc trên bảng màu, sau đó nhấn \textit{Lưu}. \newline
\quad \textit{Tự động:} Hệ thống kích hoạt logic đa ngưỡng dựa trên cảm biến ánh sáng: \newline
\qquad -- Ánh sáng > 40\%: Tắt đèn. \newline
\qquad -- Ánh sáng 10--40\%: Đèn trắng, 80\% độ sáng. \newline
\qquad -- Ánh sáng < 10\%: Đèn vàng, 30\% độ sáng. \newline
3. Lệnh được gửi qua MQTT đến mạch Yolo:bit. \newline
4. Hệ thống cập nhật trạng thái và ghi nhật ký. \\ \hline
\textbf{Luồng ngoại lệ}     & 
\textit{E1 -- Thiết bị offline:} Hệ thống thông báo không thể kết nối, lệnh không được gửi. \newline
\textit{E2 -- Ghi đè cấu hình:} Nếu thiết bị đang chạy lịch trình, chế độ thủ công sẽ ghi đè lên. \\ \hline
\textbf{Hậu điều kiện}      & Thiết bị hoạt động theo đúng chế độ; trạng thái được lưu vào system log. \\ \hline
\end{tabular}
\end{table}

% --- UC-04 ---
\paragraph{UC-04: Thiết lập lịch trình}

\begin{table}[H]
\centering
\renewcommand{\arraystretch}{1.4}
\begin{tabular}{|p{3.8cm}|p{10cm}|}
\hline
\textbf{Tên Use-case}       & Thiết lập lịch trình \\ \hline
\textbf{Mã Use-case}        & UC-04 \\ \hline
\textbf{Tác nhân}           & Người dùng \\ \hline
\textbf{Tiền điều kiện}     & Người dùng đã đăng nhập. Hệ thống và thiết bị đang hoạt động bình thường. \\ \hline
\textbf{Luồng sự kiện chính} & 
1. Người dùng truy cập mục \textit{Lập lịch thiết bị}. \newline
2. Hệ thống hiển thị danh sách lịch trình hiện có. \newline
3. Người dùng chọn \textit{Thêm mới} hoặc chỉnh sửa lịch trình cũ. \newline
4. Người dùng chọn thiết bị, ngày trong tuần, giờ bật và giờ tắt. \newline
5. Người dùng thiết lập thêm màu sắc và độ sáng (nếu là đèn RGB). \newline
6. Người dùng nhấn \textit{Lưu} và xác nhận. \newline
7. Hệ thống lưu lịch trình vào database và kích hoạt bộ đếm thời gian. \\ \hline
\textbf{Luồng ngoại lệ}     & 
\textit{E1 -- Xung đột lịch trình:} Hai lịch trình trùng thiết bị và thời điểm -- hệ thống cảnh báo yêu cầu điều chỉnh. \newline
\textit{E2 -- Ghi đè bởi thủ công:} Lịch trình sẽ bị tạm ghi đè nếu người dùng thao tác thủ công sau đó. \\ \hline
\textbf{Hậu điều kiện}      & Lịch trình được kích hoạt và thực thi đúng giờ; hành động được ghi vào system log. \\ \hline
\end{tabular}
\end{table}

% --- UC-05 ---
\paragraph{UC-05: Thiết lập cảnh báo}

\begin{table}[H]
\centering
\renewcommand{\arraystretch}{1.4}
\begin{tabular}{|p{3.8cm}|p{10cm}|}
\hline
\textbf{Tên Use-case}       & Thiết lập cảnh báo \\ \hline
\textbf{Mã Use-case}        & UC-05 \\ \hline
\textbf{Tác nhân}           & Người dùng, Hệ thống, Cảm biến \\ \hline
\textbf{Tiền điều kiện}     & Thiết bị cảm biến đang gửi dữ liệu về server định kỳ. \\ \hline
\textbf{Luồng sự kiện chính} & 
1. Người dùng truy cập mục \textit{Thiết lập cảnh báo}. \newline
2. Người dùng chọn cảm biến cần giám sát và nhập giá trị ngưỡng an toàn. \newline
3. Người dùng nhấn \textit{Lưu} và xác nhận. \newline
4. Hệ thống liên tục nhận dữ liệu từ vi điều khiển và so sánh với ngưỡng. \newline
5. Nếu giá trị vi phạm ngưỡng liên tục qua 3 chu kỳ gửi, hệ thống xác nhận cảnh báo. \newline
6. Hệ thống gửi thông báo đẩy đến điện thoại người dùng. \newline
7. Người dùng nhấn vào thông báo để xem chi tiết: vị trí, giá trị vi phạm và thời gian. \\ \hline
\textbf{Luồng ngoại lệ}     & 
\textit{E1 -- Giá trị vi phạm nhất thời:} Nếu dữ liệu chỉ vi phạm 1--2 chu kỳ rồi trở về bình thường, hệ thống không gửi cảnh báo. \newline
\textit{E2 -- Cảm biến ngắt kết nối:} Hệ thống gửi thông báo riêng về tình trạng mất dữ liệu cảm biến. \\ \hline
\textbf{Hậu điều kiện}      & Ngưỡng cảnh báo được cập nhật; cảnh báo và thao tác được ghi vào system log. \\ \hline
\end{tabular}
\end{table}

% --- UC-06 ---
\paragraph{UC-06: Nhận dữ liệu cảm biến}

\begin{table}[H]
\centering
\renewcommand{\arraystretch}{1.4}
\begin{tabular}{|p{3.8cm}|p{10cm}|}
\hline
\textbf{Tên Use-case}       & Nhận dữ liệu cảm biến \\ \hline
\textbf{Mã Use-case}        & UC-06 \\ \hline
\textbf{Tác nhân}           & Hệ thống, Vi điều khiển Yolo:bit \\ \hline
\textbf{Tiền điều kiện}     & Mạch Yolo:bit đã kết nối WiFi và xác thực thành công với Adafruit IO qua MQTT. \\ \hline
\textbf{Luồng sự kiện chính} & 
1. Mạch Yolo:bit đọc giá trị từ cảm biến nhiệt độ, độ ẩm và ánh sáng theo chu kỳ $\leq$ 5 giây. \newline
2. Vi điều khiển chuyển đổi dữ liệu thô thành giá trị vật lý và publish lên feed tương ứng trên Adafruit IO. \newline
3. Đồng thời, hiển thị giá trị tức thời lên màn hình LCD tại phòng. \newline
4. Backend ứng dụng subscribe feed MQTT, nhận dữ liệu và lưu vào database. \newline
5. App cập nhật biểu đồ và giá trị trên Dashboard với độ trễ $\leq$ 3 giây. \\ \hline
\textbf{Luồng ngoại lệ}     & 
\textit{E1 -- Mất kết nối MQTT:} Hệ thống tự động thử kết nối lại trong vòng 10 giây (NFR-R2). \newline
\textit{E2 -- Mất WiFi tại mạch:} Vi điều khiển vẫn đọc cảm biến và điều khiển thiết bị cục bộ theo ngưỡng đã lưu (NFR-O1). \\ \hline
\textbf{Hậu điều kiện}      & Dữ liệu cảm biến được lưu vào database và hiển thị trên App theo thời gian thực. \\ \hline
\end{tabular}
\end{table}

% --- UC-07 ---
\paragraph{UC-07: Xem nhật ký và báo cáo}

\begin{table}[H]
\centering
\renewcommand{\arraystretch}{1.4}
\begin{tabular}{|p{3.8cm}|p{10cm}|}
\hline
\textbf{Tên Use-case}       & Xem nhật ký và báo cáo \\ \hline
\textbf{Mã Use-case}        & UC-07 \\ \hline
\textbf{Tác nhân}           & Người dùng, Hệ thống \\ \hline
\textbf{Tiền điều kiện}     & Người dùng đã đăng nhập. Hệ thống đã lưu các hành động trong system log. \\ \hline
\textbf{Luồng sự kiện chính} & 
1. Người dùng truy cập mục \textit{Giám sát và thống kê}. \newline
2. Hệ thống hiển thị biểu đồ trực quan cảm biến theo thời gian thực. \newline
3. Người dùng chọn \textit{Nhật ký hoạt động}; hệ thống hiển thị danh sách gồm thao tác người dùng, lệnh tự động và cảnh báo. \newline
4. Người dùng lọc nhật ký theo khoảng thời gian hoặc theo phòng. \newline
5. Người dùng chọn \textit{Xuất báo cáo}, chọn khoảng thời gian thống kê. \newline
6. Hệ thống tổng hợp dữ liệu và tạo file báo cáo. \newline
7. Người dùng xác nhận tải file về thiết bị. \\ \hline
\textbf{Luồng ngoại lệ}     & 
\textit{E1 -- Không có dữ liệu trong khoảng thời gian:} Hệ thống thông báo không có nhật ký phù hợp. \newline
\textit{E2 -- Lỗi xuất file:} Hệ thống thông báo lỗi và giữ nguyên màn hình để người dùng thử lại. \\ \hline
\textbf{Hậu điều kiện}      & File báo cáo được tải về thành công; hành động xuất báo cáo được ghi vào system log. \\ \hline
\end{tabular}
\end{table}

%=============================================================
\subsection{Yêu cầu hệ thống (System Requirements)}
%=============================================================

Yêu cầu hệ thống chuyển hóa các mong muốn của người dùng thành các đặc tả kỹ thuật cụ thể cho cả phần cứng (Yolo:bit) và phần mềm (Mobile App, Server Adafruit IO).

\subsubsection{Yêu cầu chức năng hệ thống (System Functional Requirements)}

Hệ thống phải thực hiện chính xác các nhóm chức năng sau để đảm bảo tính nhất quán giữa ứng dụng di động và thiết bị thực tế:

\begin{enumerate}[label=\textbf{SR\arabic*:}]
    \item \textbf{Phân cấp quản lý dữ liệu:}
    \begin{itemize}
        \item Hệ thống phải tổ chức dữ liệu theo mô hình cây: Tài khoản $\rightarrow$ Nhà $\rightarrow$ Tầng $\rightarrow$ Phòng $\rightarrow$ Thiết bị.
        \item Khi một tầng bị xóa, tất cả các phòng và thiết bị thuộc tầng đó phải được xử lý tương ứng trong cơ sở dữ liệu (cascade delete).
    \end{itemize}

    \item \textbf{Giao diện truyền tin (MQTT Interface):}
    \begin{itemize}
        \item Hệ thống phải duy trì kết nối liên tục với Adafruit IO qua giao thức MQTT.
        \item Tự động ánh xạ (mapping) các Feed trên Adafruit với đúng ID của từng phòng trên ứng dụng Mobile.
        \item Khi kết nối bị ngắt, cơ chế reconnect phải được kích hoạt tự động trong vòng 10 giây.
    \end{itemize}

    \item \textbf{Xử lý dữ liệu cảm biến:}
    \begin{itemize}
        \item Đọc giá trị Analog/Digital từ cảm biến nhiệt độ, độ ẩm và ánh sáng từ mạch Yolo:bit.
        \item Chuyển đổi dữ liệu thô thành giá trị vật lý và gửi lên Server định kỳ $\leq$ 5 giây/lần.
        \item Hiển thị giá trị tức thời lên màn hình LCD tại phòng.
    \end{itemize}

    \item \textbf{Cơ chế điều khiển thiết bị:}
    \begin{itemize}
        \item \textbf{Manual Control:} Tiếp nhận lệnh từ App (JSON format qua MQTT), giải mã và điều khiển chân tín hiệu của Quạt hoặc Đèn RGB.
        \item \textbf{Auto Control:} Hệ thống tự động so sánh giá trị cảm biến với ngưỡng (threshold) lưu trong bộ nhớ để kích hoạt thiết bị đầu ra mà không cần lệnh từ App.
        \item \textbf{Offline Control:} Khi mất kết nối Internet, vi điều khiển sử dụng ngưỡng đã lưu cục bộ gần nhất để tiếp tục điều khiển tự động.
    \end{itemize}

    \item \textbf{Công cụ lập lịch (Scheduler Engine):}
    \begin{itemize}
        \item Bộ đếm thời gian (Real-time Clock) trên hệ thống phải đồng bộ với giờ máy chủ.
        \item Kiểm tra danh sách lịch trình (Schedule List) mỗi phút một lần để thực thi các lệnh Bật/Tắt đã cài đặt với sai số $\leq$ 5 giây.
    \end{itemize}

    \item \textbf{Hệ thống ghi nhật ký (Logging System):}
    \begin{itemize}
        \item Mọi thay đổi trạng thái (do người dùng bấm hoặc do tự động) phải được ghi lại với các trường: Thời gian, Thiết bị, Hành động, Nguồn kích hoạt (User/Auto).
        \item Toàn bộ dữ liệu cảnh báo và lịch sử thao tác phải được lưu trữ thành công vào Database.
    \end{itemize}
\end{enumerate}

\subsubsection{Giao diện hệ thống (System Interfaces)}

Hệ thống bao gồm 3 lớp giao diện tương tác chính:
\begin{itemize}
    \item \textbf{Giao diện người dùng (UI):} Ứng dụng Mobile cung cấp các Dashboard trực quan, nút bấm điều khiển và bảng thiết lập lịch trình.
    \item \textbf{Giao diện lập trình (API/Protocol):} Sử dụng RESTful API cho các tác vụ quản lý (thêm tầng/phòng) và MQTT cho các tác vụ điều khiển thời gian thực.
    \item \textbf{Giao diện phần cứng (Hardware Interface):} Các chân I/O trên mạch Yolo:bit kết nối với cảm biến và actuators, tuân thủ đúng sơ đồ nối dây của bộ kit OhStem.
\end{itemize}